% Adapted from class/easy/linkify.tex for the dataclass Link.
\begin{blocksection}
\question Write a function \texttt{linkify} that, given a list \texttt{lst}, returns a linked list with the same values. You may assume that there will be no nested lists. \\

Try writing this function both recursively and iteratively. \\

\begin{lstlisting}
def linkify_rec[T](lst: list[T]) -> LinkedList[T]:
    """
    >>> lst = [0, 1, 2, 3]
    >>> print(linkify_rec(lst))
    (0 1 2 3)
    >>> linkify_rec([])
    ()
    """
    if ___________________________________:

        __________________________________

    else:

        __________________________________


def linkify_iter[T](lst: list[T]) -> LinkedList[T]:

    ______________________________________

    for __________________________________:

        __________________________________

    ______________________________________

\end{lstlisting}
\begin{solution}[1in]
\begin{lstlisting}
def linkify_rec[T](lst: list[T]) -> LinkedList[T]:
    if not lst:
        return ()
    else:
        return Link(lst[0], linkify_rec(lst[1:]))

def linkify_iter[T](lst: list[T]) -> LinkedList[T]:
    retVal: LinkedList[T] = ()
    for elem in lst[::-1]:
        retVal = Link(elem, retVal)
    return retVal
\end{lstlisting}
\end{solution}
\end{blocksection}

\begin{questionmeta}
\textbf{Teaching Tips}\par\nopagebreak[4]
Suggested Time: 9 min recursive; add 5 min iterative; Difficulty: Medium
\nopagebreak[4]
\begin{itemize}
    \item \textbf{Read the types:} The type variable \lstinline{T} represents
    the input element type, which is preserved in the output linked list.
    \item \textbf{Use the recursive leap of faith:} Identify the base case,
    smaller input, and how to build the answer. The recursive call returns
    a linked list for the remaining Python list; empty input returns
    \lstinline{()}.
    \item \textbf{Build iteratively:} Process the input from right to left,
    placing each element at the front of the result. This preserves the
    input order. Use the recursive version first and the iterative version
    as additional practice.
\end{itemize}
\end{questionmeta}
